\chapter*[(Yellow Card)]{Second Warning to McDowell}


Needless to say, this detail has gone completely overlooked by Christian commentators. Here's what Gladkov says, for example:

\begin{smallquote}
    Most of all, [the high priests] needed to ascertain whether Jesus' body had been stolen the previous night, from Friday into Saturday, or else they would have had no reason to set a guard [...]. And they certainly ordered the stone rolled aside, verified that the Lord's body had not been stolen, and only then rolled back the stone and placed the seal upon it...
\end{smallquote}
While it doesn't directly follow from the Gospel text (``So they went and made the tomb secure, sealing the stone and setting the guard''---Mt 27:66), I actually agree with Gladkov here: there was certainly an inspection. However, let's try to imagine how this would have played out in reality.

The soldiers roll aside the stone covering the entrance to the burial chamber. The ``delegates'' from the Sanhedrin---devout Jews---are in a rather unenviable position. It's bad enough they've already defiled themselves by breaking the Sabbath; even worse, now they need to look over a tomb (recall the Jews' shock at Jesus' desire to see the deceased Lazarus). Do you really think they'd risk adding another terrifying breach of Mosaic Law---coming in contact with a dead body---to their list of Sabbath labors, just so they can check whether the tomb's occupant is actually a corpse or... \empty{an empty cocoon of burial shrouds?} I'd bet my life that at most, they stuck the very tips of their noses into the vault, then made a mad dash for the nearest synagogue to atone for their sins, yelling over their shoulder, ``Seal it up!'' The legionnaires presumably exchanged glances, twirled their fingers by their temple, then sent one of the rookies to the next village over for some hooch and leisurely started going about their duties.

So here's the version you've all been waiting for. After waiting for the women to leave, Joseph and Nicodemus extract the body from the tomb after nightfall; in its place, they leave a dummy they've prepared in advance---a fabric ``cocoon.'' In the morning, the two of them, by way of proxy, sell their colleagues in the Sanhedrin on the idea that they need to guard the tomb; they have perfectly anticipated what circumstances might prompt them to take on such a ``rush job'' on the occasion of the Passover Sabbath. The appearance of the guards at the tomb is a crucial element of the plan, which allows them to kill two birds with one stone. First, it will preemptively invalidate the high priests' inevitable attempt to write off the body's disappearance as a scheme on the part of the disciples; second, it will give the perpetrators of the hoax a convincing alibi and divert all suspicion away from them.

On the night from Saturday into Sunday, Nicodemus and Joseph---again by way of proxy---inform the procurator that his men are guarding an ``empty shell.'' An officer is immediately sent to the site, who confirms that there really is nothing inside the tomb but rags. Pilate's first thought is a perfectly natural one: he has fallen into a trap the high priests have laid for him, with the goal of compromising the procurator the Jews so strongly hate. All he can do now is go on the offensive. On his instructions, the legionnaires appear at the Sanhedrin at dawn and make a ruckus, threatening to blow the lid off this whole Jew racket and show the rudely awakened high priests for the crooks they really are. The high priests, no less embarrassed than Pilate, vow and swear that they have nothing to do with the affair, and come to a mutually satisfactory conclusion: assume that Christ's body was stolen by his disciples. After holding out a short while for appearance's sake, the soldiers agree; the procurator's order has been fulfilled, both sides have more or less saved face, and a little money has come trickling in from the Sanhedrin, and so they return to their barracks with ``a sense of pride and accomplishment.''

Meanwhile, Joseph and Nicodemus brilliantly finish up their scheme. In the most dangerous scenario, the Romans conspire with the high priests to reseal the tomb at once and make it look like nothing has happened (an act Pilate wouldn't dare carry out himself, fearing it could be a clever provocation on the part of the Sanhedrin). In that case, of course, Joseph could open the sealed tomb in front of witnesses the next day---he owned it, after all!---and ``discover'' the disappearance of the body; the problem there, however, would be that his part in the scheme would become so obvious that it could no longer be denied. For that reason, the first person to discover the open and empty tomb needs to be someone other than him. And that's where the Myrrhbearers come in.

For almost two millennia, people have read this scene over and over and surprisingly failed to ask themselves a simple question: what goal, strictly speaking, did Mary Magdalene, Mary of Jacob, and Salome have when visiting Christ's tomb at the crack of dawn? Just what rituals were they planning to perform, if they themselves had witnessed his interment on Friday and knew that their teacher had been given a ``first-class'' burial (Jn 19:39-40)? Why did they bring spices with them (hence the title ``Myrrhbearers''), if Nicodemus had already applied ``a mixture of myrrh and aloes, about a hundred pounds'' (Jn 19:39), a fact they were more than well aware of? And just where did they get the brilliant idea to \emph{disturb the peace of a dead man}---a sin in any set of human laws, and completely unthinkable in Jewish law? But if we assume that the true goal of the Myrrhbearers' pre-dawn vigil (as well as that of Peter and John, who were soon mobilized by Mary Magdalene) was to serve as witnesses to the body's disappearance and prevent the authorities from quietly resealing the tomb, then everything immediately falls into place.

So, Joseph and Nicodemus' cunning, yet risky plan went off without a hitch. But when you really think about it, what were they actually risking? The only true weak spot in their plan was the sealing of the tomb, when the body's disappearance might have theoretically been discovered---but so what? On the night from Friday into Saturday, just about anyone could have stolen the body from the unguarded grave. Even then, for what it's worth, Joseph and Nicodemus' ``indirect alibi'' would still be of some use to them; after all, in the event of an investigation, it could be shown that they were the ones who originally came up with the idea of guarding the burial site. So even in the worst possible scenario, the guys would simply ``break even.''

But what goal did they have in mind when hatching this scheme? Well, the answer to that is pretty clear. It's plainly obvious that if a major scandal were to break out, it would be the powerful high priests---Caiaphas and his father-in-law Annas---who would bear the brunt of the outrage. In the year AD 36, Caiaphas was disgracefully ``dismissed from his present position'' and even stripped of the title of high priest, which, as a rule, lasted for life. And who's to say it wasn't the scandalous story of some Galilean pretender's body disappearing almost two years previously (or alternatively: killing the Son of God) that wound up being the last straw? Moreover, Caiaphas' fall led to the rival Sadducee faction coming into power, under the leadership of the new high priest Jonathan; it was evidently this faction to which Joseph and Nicodemus belonged as well. Jonathan's policy was more pro-Roman than his predecessor's, for which he was killed by Jewish terrorists known as the Sicarii; as for Joseph, let us recall that he was a ``\emph{friend of Pilate} and the Lord.'' Looks like everything's coming together...

A simpler and more straightforward option is also possible: the stolen body was simply used as blackmail against the ruling Judean clan. ``Give us 45 gold talents/the post of deputy chief of staff of the secret police/tax farming payments in the province of Edom (delete as appropriate, fill in as necessary) and we'll return the body. Or else youse gonna regret it...'' Clearly, the deal fell through; the outcome is apparent.

From this point on, I'll let McDowell try and sort out this version. Once again, I will invoke my ``right of veto'' and simply refuse to examine this hypothesis, as it runs counter to my ``presumption of honesty''---an option that my absentee opponent, unfortunately for him, lacks entirely. As I'm sure you remember, I set out to construct a version based entirely on the integrity of the Savior's associates, and I have no intention of backing down. Of course, this makes for a more challenging task, but accordingly a far more interesting one.

As for McDowell, his position seems unenviable indeed. I have honestly tested this version for ``fracture and tension'' and have yet to find any noticeable weak spots. And considering that last time (with the ``Uncrucified Christ'' hypothesis), McDowell's ship ran into a floating mine, now it seems the captain has failed to spot an enemy submarine, which has given him a salvo from all six torpedo tubes.